% B's two content frames; F may insert them into the shared Beamer deck.
% Updated 2026-10-20 after F's actual 30-question S0 answers (PR #42).
% This is not an E scoring result or an S2 comparison.
\section{S0 普通检索基线}
\presenter{YaoYang-123456}
\source{B: reports/s0-local-run-20261019.md; F: PR \#42, s0-answer-summary-20261019.json}
\begin{frame}{S0：30题已完成真实检索与回答}
\small
\textbf{流程：}原始对话分块 $\rightarrow$ MiniLM 向量 $\rightarrow$
top-60 检索 $\rightarrow$ F 统一回答 $\rightarrow$ E 统一评分
\par\vspace{3mm}
\begin{tabularx}{\linewidth}{@{}p{33mm}Y@{}}
\toprule
固定数据 & 3个开发会话，30题（类别8/8/7/7）；前5题为联调子集\\[1mm]
B 的真实检索 & 384维、归一化、CPU；5/5和30/30题完成，每题60块；保留原始ID、时间和说话人\\[1mm]
F 的真实回答 & 5＋25次调用，30/30条生成成功；0次失败、截断或重复\\[1mm]
质量与验证 & A/B/F共29项针对性测试通过；E尚未判分，不能报告正确率\\\bottomrule
\end{tabularx}
\aside{S0是原始对话向量RAG，不是论文的平铺事实记忆消融；生成成功不等于答对。}
\note{建议约45秒。我们用普通向量检索作为StructMem的系统级基线。A冻结30题，B从原始会话检索，每题60块；来源、时间和说话人可追踪。F随后用统一模型按前5题、后25题完成30个真实回答，没有重复收费。但E还未评分，所以不能报告准确率。}
\end{frame}

\source{B: conv-42/qa\_index=15 真实检索；F: PR \#42 的真实预测与usage，尚未评分}
\begin{frame}{实例：检索证据、干扰与未完成的判分}
\small
\textbf{问题：}When did Nate get purple hair?
\par\vspace{2mm}
\begin{tabularx}{\linewidth}{@{}p{18mm}p{24mm}Y@{}}
\toprule
检索排名 & 原始消息 & 实际取回内容\\\midrule
1 & 4月15日 / D7:1 & “Dyed my hair last week”\\
2 & 同会话 / D7:3 & 图片说明包含紫色头发\\
3 & 另一会话 / D10:4 & 相似的紫发图片描述，可能干扰\\\bottomrule
\end{tabularx}
\par\vspace{3mm}
\textbf{F 的实际预测：}染发时间是2022年4月15日之前一周。\\
\textbf{生成用量：}30题输入132,531 tokens，输出421 tokens；
估算费用约\key{0.27元}（未对账）。
\par\vspace{2mm}
\aside{相对时间需要结合消息日期。此答案仍须由E按冻结标准答案评分；准确率暂无，S2尚未完整运行。}
\note{建议约45秒。这个案例在检索中同时出现时间证据、图片描述和相似干扰。F的真实回答尝试将last week绑定4月15日，但我们没有根据标准答案给它判对。全批30次模型生成的费用估算约0.27元；E负责真正的正确率和至少四个案例人工核查。B不重跑、也不代替E评分。}
\end{frame}
